% !TeX program = lualatex
% !TeX encoding = utf8
% !TeX spellcheck = uk_UA
% !BIB program = biber

\documentclass{PySCFBook}


\pagestyle{main}
\begin{document}

\section{Статичні молекулярні властивості}

У квантовій хімії статичні молекулярні властивості відіграють ключову роль у розумінні поведінки молекул під впливом зовнішніх факторів. Вони являють собою похідні енергії молекули відносно параметрів пертурбації, таких як положення ядер чи зовнішні поля. Ці властивості дозволяють прогнозувати, як молекула реагує на зміни в своєму оточенні, і є основою для багатьох розрахункових методів у хімії та фізиці. Нижче ми розглянемо визначення цих властивостей, наведемо приклади, порівняємо методи їх обчислення та обговоримо вирази з теорії пертурбацій.

\subsection{Визначення статичних молекулярних властивостей}

Розглянемо ізольовану молекулу, гамільтоніан якої \(\hat{H}_0 = \hat{T} + \hat{V}_{ne} + \hat{W}_{ee}\) описує кінетичну енергію електронів, взаємодію електронів з ядрами та взаємодію електронів між собою. Якщо до цієї системи застосувати статичну (незалежну від часу) пертурбацію \(\hat{V}(x)\), то повний гамільтоніан стає:

\[
\hat{H}(x) = \hat{H}_0 + \hat{V}(x),
\]

де параметр \(x\) визначає силу пертурбації. При \(x = 0\) пертурбація зникає, тобто \(\hat{V}(x = 0) = 0\). Ця пертурбація може моделювати різні зовнішні впливи, такі як електричні чи магнітні поля, або зміни геометрії молекули.

Енергія пертурбованої системи для певного стану \(\Psi(x)\) (зазвичай основного стану) може бути виражена як:

\[
E(x) = \frac{\langle \Psi(x) | \hat{H}(x) | \Psi(x) \rangle}{\langle \Psi(x) | \Psi(x) \rangle}.
\]

Розкладаючи цю енергію в ряд Тейлора відносно \(x\) навколо точки \(x = 0\):

\[
E(x) = E^{(0)} + E^{(1)} x + \frac{1}{2} E^{(2)} x^2 + \cdots,
\]

де \(E^{(0)}\) — енергія непертурбованої системи, а коефіцієнти \(E^{(1)}, E^{(2)}, \ldots\) є похідними енергії:

\[
E^{(1)} = \left. \frac{dE}{dx} \right|_{x=0}, \quad E^{(2)} = \left. \frac{d^2 E}{dx^2} \right|_{x=0}
\]

і так далі. Ці похідні називаються статичними молекулярними властивостями, оскільки вони описують реакцію молекули на статичні пертурбації. Вони є фундаментальними для прогнозування фізичних і хімічних характеристик молекул, таких як стабільність структур чи оптичні властивості.

\subsection{Приклади статичних молекулярних властивостей}

Статичні молекулярні властивості охоплюють широкий спектр характеристик молекул. Ось кілька ключових прикладів, які ілюструють їх застосування.

Спочатку розглянемо геометричні похідні. Тут параметр пертурбації \(x\) відповідає змінам положень ядер \(R_i - R_i^0\) у наближенні Борна-Оппенгеймера, де електронний рух відокремлений від ядерного. Молекулярний градієнт і гессіан визначаються як:

\[
\frac{\partial E}{\partial R_i}, \quad \frac{\partial^2 E}{\partial R_i \partial R_j}.
\]

Градієнт вказує на сили, що діють на ядра, і використовується для оптимізації геометрії молекули, тобто знаходження рівноважної конфігурації з мінімальною енергією. Гессіан, у свою чергу, описує кривизну поверхні потенціальної енергії і необхідний для розрахунку гармонічних вібраційних частот, які порівнюють з експериментальними спектроскопічними даними для ідентифікації молекул.

Наступний приклад — електричні властивості. Тут \(x\) — компоненти зовнішнього статичного електричного поля \(E_i\). Електричний дипольний момент і поляризовність задаються виразами:

\[
\left. \frac{\partial E}{\partial E_i} \right|_{E=0} = -\mu_i, \quad \left. \frac{\partial^2 E}{\partial E_i \partial E_j} \right|_{E=0} = -\alpha_{i,j}.
\]

Дипольний момент \(\mu_i\) характеризує асиметрію розподілу заряду в молекулі, що впливає на її взаємодію з іншими молекулами, наприклад, у розчинах. Поляризовність \(\alpha_{i,j}\) показує, наскільки легко деформувати електронну хмару під дією поля, і є ключовою для розуміння оптичних явищ, таких як рефракція світла.

Ще один важливий клас — магнітні властивості в ядерному магнітному резонансі (ЯМР). Параметри пертурбації — ядерні магнітні дипольні моменти \(m_a\) та зовнішнє магнітне поле \(B\). Константи екранування ЯМР і спін-спінової взаємодії:

\[
\frac{\partial^2 E}{\partial m_a \partial B} = -1 + \sigma_a, \quad \frac{\partial^2 E}{\partial m_a \partial m_b} = J_{a,b}.
\]

Константа екранування \(\sigma_a\) описує, як електронна хмара послаблює зовнішнє поле для ядра \(a\), що визначає хімічний зсув у спектрі ЯМР. Константа \(J_{a,b}\) характеризує взаємодію між спінами ядер через електронну систему, дозволяючи визначати структуру молекули. Аналогічно, для електронного парамагнітного резонансу (ЕПР) визначаються константи гіпердрібної взаємодії та g-фактори, які важливі для вивчення радикалів і перехідних металів.

Ці приклади демонструють, як статичні властивості пов'язують квантово-механічні розрахунки з експериментальними спостереженнями, роблячи теорію пертурбацій потужним інструментом у сучасній хімії.

\subsection{Числове проти аналітичного диференціювання}

Для обчислення похідних енергії, таких як статичні молекулярні властивості, існують два основні підходи: числове та аналітичне диференціювання. Кожен з них має свої переваги та обмеження, і вибір залежить від складності системи та вимог до точності.

Числове диференціювання базується на наближенні похідних за допомогою кінцевих різниць або поліноміальної інтерполяції. Наприклад, перша похідна може бути апроксимована як \((E(x + \Delta x) - E(x - \Delta x))/(2 \Delta x)\). Цей метод простий у реалізації, оскільки не вимагає модифікації базового коду для розрахунку енергії. Однак він має суттєві недоліки: числова точність обмежена через помилки округлення та вибору кроку \(\Delta x\), обчислювальна ефективність низька через необхідність багатьох оцінок енергії, і метод не є універсальним — зазвичай обмежується реальними статичними пертурбаціями, без урахування комплексних чи динамічних випадків.

На противагу цьому, аналітичне диференціювання передбачає явний розрахунок похідних з аналітичних виразів, отриманих диференціюванням рівнянь квантової механіки. Хоча реалізація складніша (потрібно виводити та програмувати формули для похідних), цей підхід забезпечує вищу точність, оскільки уникає наближень, і кращу обчислювальну ефективність, особливо для вищих похідних. Крім того, він більш універсальний, дозволяючи розглядати комплексні та залежні від часу пертурбації, що відкриває шлях до вивчення динамічних властивостей.

Загалом, аналітичне диференціювання є переважним для серйозних розрахунків у квантовій хімії, оскільки воно поєднує точність з ефективністю, хоча й вимагає глибшого теоретичного розуміння.

\subsection{Вираз суми по станах з теорії пертурбацій}

Для "точних" хвильових функцій, таких як повна взаємодія конфігурацій (full CI) у скінченному базисі, можна застосовувати стандартну теорію пертурбацій для виведення виразів похідних енергії. Це дає суму по станах (sum-over-states, SOS), яка є елегантним способом виразити властивості через внесок збуджених станів.

Наприклад, для лінійної пертурбації \(\hat{H}(x) = \hat{H}_0 + x \hat{V}\), перша похідна енергії:

\[
E^{(1)} = \langle \Psi_0 | \hat{V} | \Psi_0 \rangle,
\]

де \(\Psi_0\) — хвильова функція основного стану. Це просто середнє значення оператора пертурбації.

Для другої похідної:

\[
E^{(2)} = 2 \sum_{n \neq 0} \frac{|\langle \Psi_0 | \hat{V} | \Psi_n \rangle|^2}{E_0 - E_n},
\]

де індекс \(n\) пробігає по всіх збуджених станах \(\Psi_n\) з енергіями \(E_n\). Цей вираз показує, як властивості другого порядку залежать від переходів до збуджених станів: чим ближче енергії збуджених станів до основного, тим більший внесок.

SOS-вирази є інтуїтивними, оскільки пов'язують властивості з спектром збуджень молекули, і корисні для інтерпретації результатів. Однак вони не застосовні безпосередньо до наближених методів, таких як Гартрі-Фок чи теорія функціоналу густини, де потрібна більш загальна теорія відгуку для врахування релаксації хвильової функції.

\subsection{Числове проти аналітичного диференціювання}

Існує два способи розрахунку похідних енергії:

Числове диференціювання: похідні розраховуються за допомогою кінцевих різниць або поліноміального підгонки.

Воно просте в реалізації.

Числова точність обмежена.

Обчислювальна ефективність низька.

Воно не універсальне. Зазвичай, можна виконувати тільки для реальних значень та статичних пертурбацій.

Аналітичне диференціювання: похідні розраховуються явно з аналітичних виразів.

Воно складне в реалізації.

Точність вища.

Обчислювальна ефективність вища.

Воно більш універсальне. Зокрема, можливі залежні від часу пертурбації.

Загалом, аналітичне диференціювання є кращим.

\subsection{Вираз суми-по-станах з теорії пертурбацій}

Для "точних" хвильових функцій (повна взаємодія конфігурацій у базисі) ми можемо використовувати пряму теорію пертурбацій для знаходження виразу похідних енергії.

Наприклад, для пертурбації, лінійної по \(x\),

\[\hat{H}(x) = \hat{H}_0 + x \hat{V}\]

перша похідна енергії:

\[E^{(1)} = \langle \Psi_0 | \hat{V} | \Psi_0 \rangle\]

а друга похідна енергії:

\[E^{(2)} = -2 \sum_{n \neq 0} \frac{\langle \Psi_0 | \hat{V} | \Psi_n \rangle \langle \Psi_n | \hat{V} | \Psi_0 \rangle}{E_n - E_0}\]

де \(\{\Psi_n\}\) та \(\{E_n\}\) — власні стани та власні значення непертурбованого гамільтоніана \(\hat{H}_0\).

Однак, ці вирази не можна застосовувати для методів на кшталт HF або KS DFT.

Тому нам потрібна більш загальна теорія відгуку.

\subsection{Загальна статична теорія відгуку}

В методі електронної структури енергія \(E(x)\) отримується шляхом оптимізації параметрів \(p = (p_1, p_2, \dots)\) в функції енергії \(E(x, p)\) для кожного фіксованого значення \(x\). Енергія таким чином:

\[E(x) = E(x, p^o(x))\]

де \(p^o(x)\) — оптимальні значення параметрів.

Зауважте, що оптимізація не обов'язково варіаційна!

Приклади:

HF та KS DFT: \(p\) — параметри орбіталей (варіаційні).

MCSCF: \(p\) — коефіцієнти конфігурацій (варіаційні) та параметри орбіталей (варіаційні).

CI: \(p\) — коефіцієнти конфігурацій (варіаційні) та параметри орбіталей (неваріаційні).

CC: \(p\) — амплітуди кластерів (неваріаційні) та параметри орбіталей (неваріаційні).

Перша похідна \(E(x)\) відносно \(x\):

\[\frac{dE(x)}{dx} = \frac{\partial E(x, p^o)}{\partial x} + \sum_i \left. \frac{\partial E(x, p)}{\partial p_i} \right|_{p=p^o} \frac{\partial p^o_i}{\partial x}\]

(явна залежність від \(x\)) (неявна залежність від \(x\)).

\(\partial p^o_i / \partial x\) називається вектором лінійного відгуку та містить інформацію про те, як змінюється хвильова функція при пертурбації системи. Його не просто розрахувати, оскільки ми не знаємо явно залежність \(p^o_i\) від \(x\).

Якщо всі параметри варіаційні, умова нульового електронного градієнта спрощує похідну:

\[\left. \frac{\partial E(x, p)}{\partial p_i} \right|_{p=p^o} = 0 \implies \frac{dE(x)}{dx} = \frac{\partial E(x, p^o)}{\partial x}\]

тобто ми не потребуємо вектор лінійного відгуку.

Якщо всі параметри варіаційні, друга похідна \(E(x)\):

\[\frac{d^2 E(x)}{dx^2} = \frac{\partial^2 E(x, p^o)}{\partial x^2} + \sum_i \left. \frac{\partial^2 E(x, p)}{\partial x \partial p_i} \right|_{p=p^o} \frac{\partial p^o_i}{\partial x}\]

Тепер нам потрібен вектор лінійного відгуку \(\partial p^o_i / \partial x\).

Щоб отримати вектор лінійного відгуку \(\partial p^o_i / \partial x\), коли всі параметри варіаційні, ми починаємо з умови стаціонарності, яка вірна для всіх \(x\):

\[\forall x, \left. \frac{\partial E(x, p)}{\partial p_i} \right|_{p=p^o} = 0\]

Таким чином, ми можемо взяти похідну цієї рівняння відносно \(x\):

\[\left. \frac{\partial^2 E(x, p)}{\partial x \partial p_i} \right|_{p=p^o} + \sum_j \left. \frac{\partial^2 E(x, p)}{\partial p_i \partial p_j} \right|_{p=p^o} \frac{\partial p^o_j}{\partial x} = 0\]

або

\[\sum_j \left. \frac{\partial^2 E(x, p)}{\partial p_i \partial p_j} \right|_{p=p^o} \frac{\partial p^o_j}{\partial x} = -\left. \frac{\partial^2 E(x, p)}{\partial x \partial p_i} \right|_{p=p^o}\]

які є рівняннями лінійного відгуку.

Це лінійна система рівнянь для вектора лінійного відгуку \(\partial p^o_i / \partial x\).

Оскільки нас цікавлять похідні енергії, оцінені при \(x = 0\), нам потрібно розрахувати непертурбовану електронну гессіану \(\partial^2 E(x = 0, p)/\partial p_i \partial p_j\).

\subsection{Лінійний відгук HF та DFT}

\subsubsection{Лінійний відгук HF: З'єднана пертурбована Гартрі-Фок (CPHF)}

Хвильова функція HF з одним детермінантом параметризується параметрами ротації орбіталей \(\kappa\):

\[|\Phi(\kappa)\rangle = e^{\hat{\kappa}(\kappa)} |\Phi\rangle\]

де \(e^{\hat{\kappa}(\kappa)}\) виконує ротації орбіталей та записується з оператором збудження орбіталей:

\[\hat{\kappa}(\kappa) = \sum_a^{\text{occ}} \sum_r^{\text{vir}} \left( \kappa_{ar} \hat{c}^\dagger_r \hat{c}_a - \kappa^*_{ar} \hat{c}^\dagger_a \hat{c}_r \right)\]

Для статичного випадку ми можемо взяти параметри ротації орбіталей з реальними значеннями, тобто \(\kappa_{ar} = \kappa^*_{ar}\).

Без пертурбації функція енергії HF:

\[E_{\text{HF}}(\kappa) = \frac{\langle \Phi(\kappa) | \hat{H}_0 | \Phi(\kappa) \rangle}{\langle \Phi(\kappa) | \Phi(\kappa) \rangle}\]

Для теорії лінійного відгуку спочатку потрібно розрахувати електронну гессіану HF:

\[\left. \frac{\partial^2 E_{\text{HF}}(\kappa)}{\partial \kappa_{ar} \partial \kappa_{bs}} \right|_{\kappa=0} = 2 (A_{ar,bs} + B_{ar,bs})\]

де

\[A_{ar,bs} = (\varepsilon_r - \varepsilon_a) \delta_{ab} \delta_{rs} + \langle rb | as \rangle - \langle rb | sa \rangle\]

та

\[B_{ar,bs} = \langle rs | ab \rangle - \langle rs | ba \rangle\]

де \(\varepsilon_k\) — енергії спін-орбіталей HF, а \(\langle rb | as \rangle\) — двоелектронні кулонівські інтеграли над спін-орбіталями HF.

Розглянемо пертурбацію статичним електричним полем \(E\). Пертурбована енергія HF:

\[E_{\text{HF}}(E, \kappa) = \frac{\langle \Phi(\kappa) | \hat{H}_0 - \hat{\mu} \cdot E | \Phi(\kappa) \rangle}{\langle \Phi(\kappa) | \Phi(\kappa) \rangle}\]

де \(\hat{\mu}\) — оператор дипольного моменту.

Електрична дипольна поляризовність CPHF:

\[\alpha^{\text{CPHF}}_{i,j} = -\left. \frac{\partial^2 E_{\text{HF}}}{\partial E_i \partial E_j} \right|_{E=0} = -\left( \frac{\partial^2 E_{\text{HF}}}{\partial E_i \partial E_j} + \sum_a^{\text{occ}} \sum_r^{\text{vir}} \left. \frac{\partial^2 E_{\text{HF}}}{\partial E_i \partial \kappa_{ar}} \right|_{\kappa=0} \frac{\partial \kappa_{ar}}{\partial E_j} \right)\]

Оскільки явна залежність \(E_{\text{HF}}(E, \kappa)\) від \(E\) лінійна, маємо \(\partial^2 E_{\text{HF}} / \partial E_i \partial E_j = 0\).

Пертурбований електронний градієнт HF:

\[\left. \frac{\partial^2 E_{\text{HF}}}{\partial E_i \partial \kappa_{ar}} \right|_{\kappa=0} = -2 \langle r | \hat{\mu}_i | a \rangle\]

де \(\langle r | \hat{\mu}_i | a \rangle\) — перехідні дипольні моментні одноелектронні інтеграли.

Поляризовність CPHF таким чином:

\[\alpha^{\text{CPHF}}_{i,j} = \sum_a^{\text{occ}} \sum_r^{\text{vir}} 2 \langle r | \hat{\mu}_i | a \rangle \frac{\partial \kappa_{ar}}{\partial E_j}\]

Вектор лінійного відгуку \(\partial \kappa_{ar} / \partial E_j\) отримується з рівнянь лінійного відгуку HF:

\[\sum_b^{\text{occ}} \sum_s^{\text{vir}} (A_{ar,bs} + B_{ar,bs}) \frac{\partial \kappa_{bs}}{\partial E_j} = \langle r | \hat{\mu}_j | a \rangle\]

У векторно-матричних позначеннях ми отримуємо остаточно:

\[\alpha^{\text{CPHF}}_{i,j} = 2 \mu_i^T (A + B)^{-1} \mu_j\]

де \(\mu_i\) — вектор з компонентами \(\mu_{i,ar} = \langle r | \hat{\mu}_i | a \rangle\), а \(A + B\) — матриця з елементами \(A_{ar,bs} + B_{ar,bs}\).

\subsubsection{Лінійний відгук KS-DFT: З'єднана пертурбована Кон-Шем (CPKS)}

Функція енергії KS-DFT:

\[E_{\text{KS-DFT}}(\kappa) = \frac{\langle \Phi(\kappa) | \hat{T} + \hat{V}_{ne} | \Phi(\kappa) \rangle}{\langle \Phi(\kappa) | \Phi(\kappa) \rangle} + E_H[\rho_{\Phi(\kappa)}] + E_{xc}[\rho_{\Phi(\kappa)}]\]

де \(E_H[\rho]\) — гартріївський функціонал густини, а \(E_{xc}[\rho]\) — функціонал обмінно-кореляційної густини.

Електронна гессіана KS-DFT має подібний вираз, як у HF:

\[\left. \frac{\partial^2 E_{\text{KS-DFT}}(\kappa)}{\partial \kappa_{ar} \partial \kappa_{bs}} \right|_{\kappa=0} = 2 (A_{ar,bs} + B_{ar,bs})\]

де

\[A_{ar,bs} = (\varepsilon_r - \varepsilon_a) \delta_{ab} \delta_{rs} + \langle rb | as \rangle + \langle rb | f_{xc} | as \rangle\]

та

\[B_{ar,bs} = \langle rs | ab \rangle + \langle rs | f_{xc} | ab \rangle\]

де \(\varepsilon_k\) — енергії спін-орбіталей KS, \(\langle rb | as \rangle\) — двоелектронні кулонівські інтеграли над спін-орбіталями KS, а \(\langle rb | f_{xc} | as \rangle\) — двоелектронні інтеграли над ядром обмінно-кореляції \(f_{xc}(r_1, r_2) = \delta^2 E_{xc}[\rho] / \delta \rho(r_1) \delta \rho(r_2)\).

Все інше ідентичне до CPHF.

\subsection{Приклад розрахунків поляризовностей}

Якщо ми нехтуємо взаємодією електрон-електрон, маємо \(A_{ar,bs} = (\varepsilon_r - \varepsilon_a) \delta_{ab} \delta_{rs}\) та \(B_{ar,bs} = 0\), і отримуємо незв'язані Гартрі-Фок (UCHF) або незв'язані Кон-Шем (UCKS) поляризовності:

\[\alpha_{i,j}^{\text{UCHF/UCKS}} = 2 \sum_a^{\text{occ}} \sum_r^{\text{vir}} \frac{\langle r | \hat{\mu}_i | a \rangle \langle a | \hat{\mu}_j | r \rangle}{\varepsilon_r - \varepsilon_a}\]

Сферично усереднені поляризовності атомів на рівнях UCHF/CPHF та UCKS/CPKS (LDA функціонал, неконтрактовані d-aug-cc-pCV5Z базисні набори):

\begin{table}[h]
\centering
\begin{tabular}{lrrrrr}
\toprule
 & UCHF & CPHF & UCKS & CPKS & Estimated exact \\
\midrule
He & 1.00 & 1.32 & 1.81 & 1.66 & 1.38 \\
Ne & 1.98 & 2.38 & 3.48 & 3.05 & 2.67 \\
Ar & 10.1 & 10.8 & 18.0 & 12.0 & 11.1 \\
Kr & 15.9 & 16.5 & 27.8 & 18.0 & 16.8 \\
Be & 30.6 & 45.6 & 80.6 & 43.8 & 37.8 \\
Mg & 55.2 & 81.6 & 122 & 71.4 & 71.3 \\
Ca & 125 & 185 & 277 & 149 & 157 \\
\bottomrule
\end{tabular}
\end{table}

Витягнуто з J. Toulouse, E. Rebolini, T. Gould, J. F. Dobson, P. Seal, and J. G. ´Angy´an, J. Chem. Phys. 138, 194106 (2013).

\section{Динамічні молекулярні властивості}

\subsection{Рівняння Шредінгера, залежне від часу, та квазі-енергія}

Тепер розглянемо періодичну залежну від часу пертурбацію:

\[\hat{H}(t) = \hat{H}_0 + x^+ \hat{V} e^{-i \omega t} + x^- \hat{V}^\dagger e^{+i \omega t}\]

де \(x^+\) та \(x^-\) представляють сили пертурбації (зрештою \(x^+ = x^-\)).

Асоційована хвильова функція \(|\bar{\Psi}(t)\rangle\) задовольняє рівняння Шредінгера, залежне від часу:

\[\left( \hat{H}(t) - i \frac{\partial}{\partial t} \right) |\bar{\Psi}(t)\rangle = 0\]

Після витягування фазового фактора \(|\bar{\Psi}(t)\rangle = e^{-i F(t)} |\Psi(t)\rangle\), його можна переформулювати як:

\[\left( \hat{H}(t) - i \frac{\partial}{\partial t} \right) |\Psi(t)\rangle = \dot{F}(t) |\Psi(t)\rangle \quad (1)\]

де \(\dot{F}(t) = dF/dt\) — залежна від часу квазі-енергія.

Квазі-енергія \(Q\) визначається як середнє за часом \(\dot{F}(t)\) за період \(T = 2\pi / \omega\):

\[Q = \frac{1}{T} \int_0^T dt \, \dot{F}(t) = \frac{1}{T} \int_0^T dt \, \frac{\langle \Psi(t) | \left( \hat{H}(t) - i \frac{\partial}{\partial t} \right) | \Psi(t) \rangle}{\langle \Psi(t) | \Psi(t) \rangle}\]

Вона задовольняє стаціонарний принцип, подібний до незалежного від часу:

\[\delta_{\Psi(t)} Q = 0 \iff \Psi(t) \text{ є розв'язком (1)}\]

\subsection{Визначення динамічних молекулярних властивостей}

Квазі-енергію можна розкласти відносно сил пертурбації \(x^+\) та \(x^-\):

\[Q = Q^{(0)} + Q^{(10)} x^+ + Q^{(01)} x^- + \frac{1}{2} Q^{(20)} (x^+)^2 + Q^{(11)} x^+ x^- + \frac{1}{2} Q^{(02)} (x^-)^2 + \cdots\]

де \(Q^{(0)} = E^{(0)}\) — непертурбована енергія, а похідні квазі-енергії:

\[Q^{(10)} = \left. \frac{\partial Q}{\partial x^+} \right|_{x=0}; \quad Q^{(01)} = \left. \frac{\partial Q}{\partial x^-} \right|_{x=0}; \quad Q^{(20)} = \left. \frac{\partial^2 Q}{\partial (x^+)^2} \right|_{x=0}; \quad Q^{(11)} = 2 \left. \frac{\partial^2 Q}{\partial x^+ \partial x^-} \right|_{x=0}; \quad \text{тощо...}\]

є \(\omega\)-залежними динамічними молекулярними властивостями.

Ми розглянемо конкретний випадок пертурбації взаємодії електричного диполя:

\[\hat{H}(t) = \hat{H}_0 - E^+ \cdot \hat{\mu} e^{-i \omega t} - E^- \cdot \hat{\mu} e^{+i \omega t}\]

де \(\hat{\mu}\) — оператор дипольного моменту, а \(E^+ = E^-\) — амплітуда осцилюючого електричного поля. Динамічна електрична дипольна поляризовність визначається як:

\[\alpha_{i,j}(\omega) = -\left. \frac{\partial^2 Q}{\partial E^-_i \partial E^+_j} \right|_{E=0}\]

Використовуючи залежну від часу теорію пертурбацій, ми отримуємо точний вираз:

\[\alpha_{i,j}(\omega) = -\left[ \sum_{n \neq 0} \frac{\langle \Psi_0 | \hat{\mu}_i | \Psi_n \rangle \langle \Psi_n | \hat{\mu}_j | \Psi_0 \rangle}{\omega - \omega_n} - \sum_{n \neq 0} \frac{\langle \Psi_0 | \hat{\mu}_j | \Psi_n \rangle \langle \Psi_n | \hat{\mu}_i | \Psi_0 \rangle}{\omega + \omega_n} \right]\]

де \(\{\Psi_n\}\) — стани \(\hat{H}_0\), а \(\{\omega_n\}\) — енергії збудження.

\subsection{Поляризовність з динамічної теорії відгуку}

У залежному від часу методі електронної структури квазі-енергія \(Q(E^+, E^-)\) отримується шляхом підставки оптимальних параметрів \(p^o(E^+, E^-)\) у функцію квазі-енергії \(Q(E^+, E^-; p)\):

\[Q(E^+, E^-) = Q(E^+, E^-; p^o(E^+, E^-))\]

На відміну від статичного випадку, параметри тепер залежать від часу та приймають комплексні значення. Вони можуть бути розкладені в моди Фур'є:

\[p = p^+ e^{-i \omega t} + p^{-*} e^{+i \omega t}\]

Ми розглянемо тільки випадок, коли всі параметри варіаційні, тобто задовольняють умови стаціонарності:

\[\left. \frac{\partial Q}{\partial p^{+*}} \right|_{p=p^o} = 0 \quad \text{та} \quad \left. \frac{\partial Q}{\partial p^{-*}} \right|_{p=p^o} = 0\]

Завдяки стаціонарності параметрів, тільки явна залежність від \(E^+, E^-\) в \(Q\) вносить вклад для властивостей першого порядку:

\[\frac{\partial Q}{\partial E^-_i} = \frac{\partial Q}{\partial E^-_i} + \left[ \frac{\partial Q}{\partial p^+} \frac{\partial p^+}{\partial E^-_i} + \frac{\partial Q}{\partial p^-} \frac{\partial p^-}{\partial E^-_i} + \text{c.c.} \right] = 0\]

Для властивостей другого порядку неявна залежність від \(E^+, E^-\) через параметри вносить вклад. Зокрема, для динамічної поляризовності ми отримуємо:

\[\alpha_{i,j}(\omega) = -\frac{\partial^2 Q}{\partial E^-_i \partial E^+_j} = -\frac{\partial^2 Q}{\partial p^+ \partial E^-_i} \frac{\partial p^+}{\partial E^+_j} - \frac{\partial^2 Q}{\partial p^- \partial E^-_i} \frac{\partial p^-}{\partial E^+_j}\]

Вектор лінійного відгуку \((\partial p^+ / \partial E^+_j, \partial p^- / \partial E^+_j)\) знаходиться з рівнянь лінійного відгуку:

\[\begin{pmatrix} \frac{\partial^2 Q}{\partial p^{+*} \partial p^+} & \frac{\partial^2 Q}{\partial p^{+*} \partial p^-} \\ \frac{\partial^2 Q}{\partial p^{-*} \partial p^+} & \frac{\partial^2 Q}{\partial p^{-*} \partial p^-} \end{pmatrix} \begin{pmatrix} \frac{\partial p^+}{\partial E^+_j} \\ \frac{\partial p^-}{\partial E^+_j} \end{pmatrix} = - \begin{pmatrix} \frac{\partial^2 Q}{\partial E^+_j \partial p^{+*}} \\ \frac{\partial^2 Q}{\partial E^+_j \partial p^{-*}} \end{pmatrix}\]

Динамічна поляризовність таким чином має загальний вираз:

\[\alpha_{i,j}(\omega) = \begin{pmatrix} \frac{\partial^2 Q}{\partial p^+ \partial E^-_i} & \frac{\partial^2 Q}{\partial p^- \partial E^-_i} \end{pmatrix}^T \begin{pmatrix} \frac{\partial^2 Q}{\partial p^{+*} \partial p^+} & \frac{\partial^2 Q}{\partial p^{+*} \partial p^-} \\ \frac{\partial^2 Q}{\partial p^{-*} \partial p^+} & \frac{\partial^2 Q}{\partial p^{-*} \partial p^-} \end{pmatrix}^{-1} \begin{pmatrix} \frac{\partial^2 Q}{\partial E^+_j \partial p^{+*}} \\ \frac{\partial^2 Q}{\partial E^+_j \partial p^{-*}} \end{pmatrix}\]

\subsection{Лінійний відгук TDHF та TDDFT}

\subsubsection{Лінійний відгук залежного від часу Гартрі-Фок (TDHF)}

Функція квазі-енергії TDHF:

\[Q_{\text{TDHF}}(E^+, E^-; \kappa) = \frac{1}{T} \int_0^T dt \, \frac{\langle \Phi(\kappa) | \hat{H}(t) - i \frac{\partial}{\partial t} | \Phi(\kappa) \rangle}{\langle \Phi(\kappa) | \Phi(\kappa) \rangle}\]

Після розрахунку всіх похідних \(Q_{\text{TDHF}}\) відносно параметрів ротації орбіталей \(\kappa_{ar}\), ми отримуємо поляризовність TDHF (для орбіталей з реальними значеннями):

\[\alpha_{i,j}^{\text{TDHF}}(\omega) = \begin{pmatrix} \mu_i \\ \mu_i \end{pmatrix}^T \left( \begin{pmatrix} A & B \\ B & A \end{pmatrix} - \omega \begin{pmatrix} 1 & 0 \\ 0 & -1 \end{pmatrix} \right)^{-1} \begin{pmatrix} \mu_j \\ \mu_j \end{pmatrix}\]

де \(A\) та \(B\) — ті самі матриці, введені в статичній теорії відгуку (тобто CPHF).

Енергії збудження TDHF \(\omega_n\) відповідають полюсам \(\alpha_{i,j}^{\text{TDHF}}(\omega)\) в \(\omega\) (тобто значенням \(\omega\), де \(\alpha_{i,j}^{\text{TDHF}}(\omega)\) розбігається). Вони знаходяться шляхом розв'язання псевдо-ермітового рівняння власних значень:

\[\begin{pmatrix} A & B \\ B & A \end{pmatrix} \begin{pmatrix} X_n \\ Y_n \end{pmatrix} = \omega_n \begin{pmatrix} 1 & 0 \\ 0 & -1 \end{pmatrix} \begin{pmatrix} X_n \\ Y_n \end{pmatrix}\]

Сили осциляторів можуть бути розраховані з власних векторів \((X_n, Y_n)\):

\[f_n = \frac{2}{3} \omega_n \sum_i \left[ (X_n + Y_n)^T \mu_i \right]^2\]

\subsubsection{Лінійний відгук залежного від часу DFT (TDDFT)}

Функція квазі-енергії TDDFT:

\[Q_{\text{TDDFT}}(E^+, E^-; \kappa) = \frac{1}{T} \int_0^T dt \left[ \frac{\langle \Phi(\kappa) | \hat{h}(t) - i \frac{\partial}{\partial t} | \Phi(\kappa) \rangle}{\langle \Phi(\kappa) | \Phi(\kappa) \rangle} + E_H[\rho_{\Phi(\kappa)}] \right] + Q_{xc}[\rho_{\Phi(\kappa)}]\]

де \(\hat{h}(t)\) — одноелектронний гамільтоніан, а \(Q_{xc}[\rho]\) — обмінно-кореляційний функціонал квазі-енергії, що залежить від залежної від часу густини \(\rho\).

На практиці ми завжди використовуємо адіабатичне наближення:

\[Q_{xc}[\rho] \approx \frac{1}{T} \int_0^T dt \, E_{xc}[\rho]\]

де \(E_{xc}[\rho]\) — обмінно-кореляційний функціонал основного стану.

Тоді ми отримуємо ті самі рівняння, як у TDHF, з тими самими матрицями \(A\) та \(B\), як у CPKS:

\[A_{ar,bs} = (\varepsilon_r - \varepsilon_a) \delta_{ab} \delta_{rs} + \langle rb | as \rangle + \langle rb | f_{xc} | as \rangle\]

\[B_{ar,bs} = \langle rs | ab \rangle + \langle rs | f_{xc} | ab \rangle\]

У точному TDDFT обмінно-кореляційне ядро \(f_{xc}\) повинно залежати від \(\omega\), але в адіабатичному наближенні воно не залежить. Це має наслідок, що подвійні або вищі електронні збудження не описуються.

\subsubsection{Наближення Тамма-Данкоффа (TDA)}

Наближення Тамма-Данкоффа (TDA) відповідає нехтуванню матрицею \(B\).

Рівняння лінійного відгуку TDHF або TDDFT тоді спрощуються до:

\[A X_n = \omega_n X_n\]

У TDHF матриця \(A\) відповідає матриці гамільтоніана (зсунутого на енергію HF) в базисі однозбуджених детермінантів:

\[A_{ar,bs} = \langle \Phi_a^r | \hat{H}_0 - E_{\text{HF}} | \Phi_b^s \rangle.\]

Таким чином, ми відновлюємо розрахунок взаємодії конфігурацій синглів (CIS).

У TDA ми отримуємо хвильові функції збуджених станів як:

\[|\Psi_n\rangle = \sum_a^{\text{occ}} \sum_r^{\text{vir}} X_{n,ar} |\Phi_a^r\rangle\]

де \(X_{n,ar}\) — коефіцієнт для одної орбітальної збудження \(a \to r\). Зазвичай, одна орбітальна збудження домінує над усіма іншими, і це використовується для "присвоєння" збудженого стану цій орбітальній збудженні.

\subsection{Приклад розрахунків енергій збудження}

Валентні та рідбергівські енергії збудження молекули CO (базисний набір Sadlej+)

\begin{table}[h]
\centering
\small
\begin{tabular}{lccrrrrr}
\toprule
Стан & Перехід & TDHF & TDHF-TDA & TDLDA & TDLDA-TDA & EOM-CCSD & Експт \\
\midrule
\multicolumn{8}{c}{Валентні енергії збудження (eV)} \\
$^3\Pi$ & $5\sigma \to 2\pi^*$ & 5.28 & 5.85 & 5.95 & 6.04 & 6.45 & 6.32 \\
$^3\Sigma^+$ & $1\pi \to 2\pi^*$ & 6.33 & 7.79 & 8.38 & 8.54 & 8.42 & 8.51 \\
$^1\Pi$ & $5\sigma \to 2\pi^*$ & 8.80 & 9.08 & 8.18 & 8.42 & 8.76 & 8.51 \\
$^3\Delta$ & $1\pi \to 2\pi^*$ & 7.87 & 8.74 & 9.16 & 9.20 & 9.39 & 9.36 \\
$^3\Sigma^-$ & $1\pi \to 2\pi^*$ & 9.37 & 9.73 & 9.84 & 9.84 & 9.97 & 9.88 \\
$^1\Sigma^-$ & $1\pi \to 2\pi^*$ & 9.37 & 9.73 & 9.84 & 9.84 & 10.19 & 9.88 \\
$^1\Delta$ & $1\pi \to 2\pi^*$ & 9.96 & 10.15 & 10.31 & 10.33 & 10.31 & 10.23 \\
$^3\Pi$ & $4\sigma \to 2\pi^*$ & 13.05 & 13.31 & 11.40 & 11.43 & 12.49 &  \\
\multicolumn{8}{c}{Рідбергівські енергії збудження (eV)} \\
$^3\Sigma^+$ & $5\sigma \to 6\sigma$ & 11.07 & 11.18 & 9.55 & 9.56 & 10.60 & 10.40 \\
$^1\Sigma^+$ & $5\sigma \to 6\sigma$ & 12.23 & 12.27 & 9.93 & 9.95 & 11.15 & 10.78 \\
$^3\Sigma^+$ & $5\sigma \to 7\sigma$ & 12.40 & 12.42 & 10.26 & 10.26 & 11.42 & 11.30 \\
$^1\Sigma^+$ & $5\sigma \to 7\sigma$ & 12.78 & 12.79 & 10.47 & 10.50 & 11.64 & 11.40 \\
$^3\Pi$ & $5\sigma \to 3\pi$ & 12.52 & 12.60 & 10.39 & 10.39 & 11.66 & 11.55 \\
$^1\Pi$ & $5\sigma \to 3\pi$ & 12.87 & 12.88 & 10.48 & 10.50 & 11.84 & 11.53 \\
\multicolumn{8}{c}{Поріг іонізації: $-\epsilon_{\text{HOMO}}$ (eV)} \\
 &  & 15.11 & 15.11 & 9.12 & 9.12 & 15.58 &  \\
\multicolumn{8}{c}{Середні абсолютні відхилення енергій збудження відносно EOM-CCSD (eV)} \\
Валентні &  & 0.89 & 0.49 & 0.37 & 0.33 &  &  \\
Рідбергівські &  & 0.93 & 0.97 & 1.21 & 1.19 &  &  \\
Всього &  & 0.91 & 0.69 & 0.73 & 0.70 &  &  \\
\bottomrule
\end{tabular}
\caption{Витягнуто з E. Rebolini, A. Savin, and J. Toulouse, Mol. Phys. 111, 1219 (2013)}
\end{table}

\subsection{Покращення наближень, використовуваних у лінійному відгуку TDDFT}

Існують поточні зусилля для покращення наближень, використовуваних у лінійному відгуку TDDFT:

Гібридні наближення (змішування HF обміну та DFT обміну) та розділені за дальністю гібридні наближення (змішування далекодійного HF обміну та короткодійного DFT обміну) покращують опис нелокальних збуджень (наприклад, рідбергівських та перенесення заряду).

Було запропоновано ряд подвійно-гібридних наближень (змішування PT2 кореляції та DFT кореляції) для TDDFT.

Підхід, подібний до TDDFT, відомий як рівняння Бете-Салпетера (BSE), спочатку з фізики конденсованого стану, все частіше застосовується до молекул. Він подібний до TDHF, але з екранованим ядром HF обміну, що враховує кореляцію.

Ситуації зі статичною/сильною кореляцією (дисоціація зв'язків, перехідні метали, конічні перетини, ...) зазвичай є проблемою в TDDFT з звичайними наближеннями. Розробляються ряд підходів, що поєднують багатоконфігураційні методи хвильових функцій (наприклад, MCSCF) з TDDFT.

\end{document}